\documentclass[12pt]{article}
\usepackage[utf8]{inputenc}
\usepackage[T1]{fontenc}
\usepackage{hyperref}
\title{Tamazight, Temacine Language Resources for K-12 Curriculum Development}
\author{Compiled by Claude Sonnet 5 (claude-sonnet-5) for Andr\'as Kornai}
\date{}
\begin{document}
\maketitle

\section{Introduction}
Temacine Tamazight (ISO 639-3 tjo, Glottocode tema1243), also called Tugurt or Oued Righ Berber, is a Zenati Berber variety of Algeria. It is spoken in the oases around Touggourt, mainly in Temacine, Blidet Amor, Meggarine and Ghomra. Wikipedia lists about 8,400 speakers (2022). The older Endangered Languages entry gives 8,000 to 10,000 and rates it as threatened. It is closely related to Mozabite (Tumzabt) and Ouargli (Teggargrent). Tamazight is a national and official language of Algeria, but one source notes that this variety has no practical official presence.

\section{Primary Educational Resources}
I found no primers, textbooks or school materials specific to Temacine Tamazight. Materials for standard Algerian Tamazight may exist, but I did not find or verify any, and they should not be assumed to match this variety.

\section{Linguistic and Reference Materials}
\begin{itemize}
\item Tugurt language (Wikipedia): \url{https://en.wikipedia.org/wiki/Tugurt_language}
\item Endangered Languages Project entry: \url{https://endangeredlanguages.com/node/94298}
\item Temacine (town): \url{https://www.wikipedia.org/wiki/Temacine}
\end{itemize}
I did not find a grammar or dictionary of this variety. Many fields in the Endangered Languages entry are blank.

\section{Religious and Community Resources}
None found.

\section{Digital and Audio Resources}
None found specific to this variety.

\section{Curriculum Development Context}
The language has several thousand speakers but little documentation online. A curriculum would need new fieldwork with community members in Temacine and neighbouring villages, since the villages differ dialectally. Algeria's constitutional recognition of Tamazight may allow some school support, but I found nothing concrete for this variety.
\end{document}